% ============================================================
%  SECTION 9 — ETHICAL CONSIDERATIONS AND REPRODUCIBILITY
%
%  Place this section immediately BEFORE \section{Conclusion}
%  in the main .tex file. For IEEE Access two-column format the
%  ordering is:
%    §1  Abstract
%    §2  Introduction
%    §3  Related Work
%    §4  Dataset Construction
%    §5  Leak Audit and Dataset Validation
%    §6  Baseline Experiments
%    §7  Discussion
%    §8  Ethical Considerations and Reproducibility  ← this file
%    §9  Conclusion
%    Appendix A  Dataset Documentation
%
%  No new \bibitem keys: all citations here were already
%  declared in related_work_and_references.tex or methodology.tex.
%  Cross-references used: sec:methodology, ssec:preprocess,
%  tab:generators, alg:preprocess, sec:experiments,
%  ssec:exp_setup, sec:datasheet.
% ============================================================

\section{Ethical Considerations and Reproducibility}
\label{sec:ethics}

\subsection{Ethical Considerations}
\label{ssec:ethics_content}

\paragraph{Intended scope and dual-use risk.}
\textbf{\datasetname} is constructed exclusively for the purpose of
training and evaluating AI-generated image detectors.
Like all resources in this space, it carries an inherent dual-use
risk: the per-generator detectability scores reported in
Section~\ref{sec:experiments} could, in principle, be used to
identify which generators are harder to detect and to preferentially
deploy those generators in adversarial settings.
We acknowledge this risk while noting that the detectability
characteristics we report are a property of the generators themselves,
not of our dataset — any sufficiently large dataset covering the same
six generators would reveal the same relative ordering.
Publishing these findings openly allows the defence community to
prioritise the development of detectors for harder-to-detect
generators rather than leaving that knowledge asymmetrically available
to adversarial actors.

\paragraph{Real image provenance and privacy.}
Real images are sourced from the MS-COCO 2017 validation
split~\cite{lin2014coco} and the ImageNet-1k validation
split~\cite{deng2009imagenet}, both of which are publicly released
research datasets with established data governance.
COCO images are sourced from Flickr under Creative Commons licences
agreed to by the photographers at upload time.
ImageNet images are distributed under ILSVRC terms that permit
non-commercial research use.
No personally identifiable information was collected by us; we do not
retain any metadata beyond what is present in the original dataset
releases.

The real images may contain depictions of people in public settings.
We did not perform any face detection, re-identification, or
demographic annotation on these images, and we do not recommend using
the real-image subset for any task involving individual identification.

\paragraph{Synthetic image safety.}
Generation pipelines were executed with
\texttt{add\_watermarker=False} and \texttt{safety\_checker=None}.
The rationale for disabling the safety checker is technical, not
permissive: enabling it would introduce a class-asymmetric
post-processing step — AI images filtered and replaced, real images
untouched — that would constitute a preprocessing branch on the class
label, precisely the kind of bias the dataset is designed to
eliminate.
We manually inspected a random sample of 200 generated images across
all six generators and found no sexually explicit, violent, or
otherwise harmful content.
The BLIP-2 captions used as prompts are derived from photographs of
everyday scenes (COCO) and object-centric images (ImageNet); they do
not describe harmful content.

\paragraph{Licence and intended use restriction.}
As detailed in Appendix~\ref{sec:datasheet}, the effective licence
for the dataset is non-commercial research use only, inherited from
the ImageNet component.
We explicitly do not authorise the use of this dataset to train
commercial detection products without first verifying compliance with
the ILSVRC terms of use.
We also do not authorise the use of detectors trained on this dataset
as the sole basis for consequential decisions (e.g., content
moderation, legal proceedings, or journalistic attribution) without
extensive additional validation beyond the benchmark conditions
reported here.

\paragraph{Misuse of detectability scores.}
The sniff-test and trained-detector results reported in
Section~\ref{sec:experiments} and Table~\ref{tab:ablation_summary} characterise the intrinsic forensic
difficulty of six generators under controlled conditions.
They should not be interpreted as absolute claims about how detectable
any given generator's output is in the wild, where post-processing
(compression, resizing, social-media reencoding) may alter artefact
profiles substantially.
A detector achieving AUROC~$= 0.9807$ on our test split should not be
assumed to achieve equivalent performance on unpreprocessed images
from the same generators.

%-------------------------------------------------------------------
\subsection{Reproducibility Statement}
\label{ssec:reproducibility}
%-------------------------------------------------------------------

We have taken the following concrete steps to ensure the complete
reproducibility of all results in this paper.

\paragraph{Dataset.}
The complete dataset — all 70,000 images, metadata, manifest, split
assignments, captions, validation report, and generation configuration
— is publicly available at\\
\url{https://huggingface.co/datasets/Shanmuk4622/ai-detection-dataset-v2}
without authentication.
The repository is tagged at version \texttt{v2.0} to ensure a stable,
citable snapshot.

\paragraph{Build pipeline.}
The complete eleven-notebook build pipeline (NB00--NB11) is included
in the repository under \texttt{notebooks/}.
Each notebook is self-contained, downloads all dependencies from
PyPI and HuggingFace, and reads its configuration exclusively from
\texttt{config.json} — no notebook requires manual parameter editing.
All random state is controlled by the single scalar \texttt{MASTER\_SEED}~$= 42$
declared in NB00; changing this value and re-running the full pipeline
produces a statistically independent dataset with identical structure.

\paragraph{Preprocessing.}
The canonical preprocessing function
(\texttt{canonical\_preprocess} in \texttt{ai\_detect\_common.py},
reproduced in Algorithm~\ref{alg:preprocess}, Section~\ref{ssec:preprocess}) is deterministic and
takes only a PIL image as input.
Running it on any image — real or AI — produces a byte-identical
output given the same input, verifiable via the SHA-256 digests
stored in the schema.

\paragraph{Generator configurations.}
All generation hyperparameters (model identifiers, number of
inference steps, classifier-free guidance scales, native resolutions,
and scheduler configurations) are recorded in \texttt{config.json}
and reproduced in Table~\ref{tab:generators}.
Seed assignments are computed deterministically by
Equation~\ref{eq:seed}; given the same \texttt{MASTER\_SEED}, model
identifier, and \texttt{source\_real\_id}, the seed is uniquely
determined and can be independently verified.

\paragraph{Baseline experiments.}
The training setup for all three baseline detectors is described
completely in Section~\ref{ssec:exp_setup}: model checkpoints
(\texttt{timm} identifiers), optimiser (AdamW,
$\text{lr} = 3 \times 10^{-4}$, weight decay $10^{-4}$),
scheduler (cosine annealing, $\eta_{\min} = 3 \times 10^{-6}$),
batch size (64 per GPU, dual T4, DataParallel),
augmentation (horizontal flip, colour jitter $\pm 0.1$),
label smoothing ($\varepsilon = 0.05$), gradient clipping (norm 1.0),
and the checkpoint selection criterion (best validation AUROC).
Trained model checkpoints are archived at\\
\pathtt{results/baseline\_experiments/checkpoints/}\\
in the same HuggingFace repository.
Per-epoch training histories are stored as JSON files
(\pathtt{resnet50\_progress.json}, \pathtt{efficientnet\_b4\_progress.json},
\pathtt{vit\_b16\_progress.json}) alongside the final
\texttt{baseline\_results.json} containing all reported metrics.

\paragraph{Hardware and software.}
All experiments were run on Kaggle GPU sessions equipped with dual
NVIDIA Tesla T4 GPUs (15.6~GB VRAM each).
The software environment is Python 3.12.13, PyTorch 2.10.0+cu128,
\texttt{timm}~0.9+, \texttt{transformers}~5.0.0,
\texttt{huggingface\_hub}~0.23+, \texttt{pyarrow}, and
\texttt{scikit-learn}.
No custom CUDA kernels or proprietary libraries are used.
All packages are installable from standard PyPI without version pinning
beyond the minimum versions listed in the notebooks.

\paragraph{Numerical results.}
All AUROC, F1, and accuracy values in Tables~\ref{tab:sniff},
\ref{tab:baselines}, \ref{tab:per_gen}, and \ref{tab:cross_gen} are
drawn directly from \texttt{baseline\_results.json} without rounding
beyond four decimal places.
The cross-generator generalisation results in Table~\ref{tab:cross_gen}
were each produced by a fresh ResNet-50 initialisation (no checkpoint
reuse across held-out runs), trained for five epochs on the reduced
training set with the same hyperparameters as the main experiment.
